% ===========================================================================
%  osc_results.tex -- result tables + per-family discussion for the
%  oscillatory suite.  \input by jsc_paper_main.tex inside
%  \section{Complex versus real networks on the oscillatory suite}.
%  Numbers are means over seeds from results/*_v3.csv (H20, cosine schedule).
% ===========================================================================

% ---------------------------------------------------------------------------
\subsection{High-frequency second-order problems}\label{ssec:res-hf}

Here the obstacle is spectral bias rather than derivative order. Table~\ref{tab:res-hf}
collects the constant-coefficient Helmholtz frequency sweep, its
variable-coefficient ``scattering'' variant, and the non-separable radial
chirp. The complex \(\sinh\) network is the most accurate method on every
constant-coefficient Helmholtz wavenumber and on every chirp frequency; on the
heterogeneous (variable-\(\kappa\)) problem it leads only at the highest
wavenumber, with a Fourier-feature or SIREN network ahead at low frequency. The
single anisotropic case \((a_1,a_2)=(1,4)\) is also lost: at this low, strongly
directional frequency a SIREN reaches \(2.2\times10^{-1}\) against the complex
network's \(3.9\times10^{-1}\), consistent with SIREN's strength on separable
low-frequency sinusoids.
The chirp is the cleanest expressivity signal: because its solution is not a
product of one-dimensional sines, the sine-based MscaleDNN -- competitive on
separable targets -- degrades sharply (errors above one), while the complex
network retains the smallest error throughout.

\begin{table}[htbp]
\centering
\small
\caption{High-frequency second-order problems: relative \(L^2\) error
\eqref{eq:relL2} (lower is better; best in bold). \(a\) is the swept wavenumber.}
\label{tab:res-hf}
\begin{tabular}{@{}llrrrr@{}}
\toprule
family & \(a\) & complex \(\sinh\) & Fourier & SIREN & MscaleDNN \\
\midrule
Helmholtz                & 2  & \textbf{5.22e-2} & 1.52e-1 & 8.09e-2 & 6.99e-1 \\
                         & 4  & \textbf{3.22e-1} & 4.32e-1 & 3.46e-1 & 1.00e0 \\
                         & 6  & \textbf{5.71e-1} & 8.51e-1 & 8.19e-1 & 1.00e0 \\
                         & 8  & \textbf{6.29e-1} & 9.32e-1 & 9.27e-1 & 1.00e0 \\
                         & 10 & \textbf{7.71e-1} & 9.86e-1 & 9.83e-1 & 1.00e0 \\
\midrule
Helmholtz (anis.\ \(1,4\)) & -- & 3.89e-1 & 7.79e-1 & \textbf{2.21e-1} & 1.00e0 \\
\midrule
Helmholtz (var.\ \(\kappa\)) & 2 & 3.59e-1 & \textbf{2.11e-1} & 3.15e-1 & 1.00e0 \\
                         & 4  & 4.64e-1 & 5.45e-1 & \textbf{4.37e-1} & 1.00e0 \\
                         & 6  & \textbf{7.12e-1} & 8.72e-1 & 8.07e-1 & 1.00e0 \\
\midrule
radial chirp             & 2  & \textbf{3.76e-1} & 7.07e-1 & 8.74e-1 & 7.93e-1 \\
                         & 4  & \textbf{4.52e-1} & 9.19e-1 & 7.16e-1 & 1.60e0 \\
                         & 6  & \textbf{8.47e-1} & 1.18e0  & 9.76e-1 & 1.89e0 \\
                         & 8  & \textbf{8.30e-1} & 1.24e0  & 9.64e-1 & 2.39e0 \\
\bottomrule
\end{tabular}
\end{table}

% ---------------------------------------------------------------------------
\subsection{High-order real operators}\label{ssec:res-ho}

Table~\ref{tab:res-ho} reports the order-controlled families. The one-dimensional
polyharmonic sweep isolates the differentiation order on a single monomial
\(\partial_x^{2m}\): the complex network is one to two orders of magnitude more
accurate than every real baseline at orders two and four, retains the lead at
order eight, and is edged out by MscaleDNN only at order six; at order ten the
fixed budget is insufficient for any method. In two dimensions the gap widens
with order, reaching a \(13\times\) advantage on the order-six triharmonic where
all real baselines have effectively failed. The anisotropic plate
(\(\Delta^2\), mixed modes \((m,m{+}1)\)) is won at every mode by a clear margin,
and the isotropic plate is won at the higher modes with an advantage that grows
from parity to \(1.85\times\). The one-dimensional Euler--Bernoulli beam is the
subtler case: a Fourier-feature network ties or wins on the smoothest modes,
because a one-dimensional fourth-order sine is precisely the stationary
single-frequency signal that Fourier features are designed for.
The exception is the third-order linearized KdV: there a SIREN network is the
most accurate at the higher wavenumbers, and the complex network leads only at
one wavenumber, reflecting that a first-order-in-time dispersive operator does
not present the balanced repeated-index structure that the complex ansatz
exploits. The nonlinear Cahn--Hilliard phase field returns to the favourable
regime: the complex network wins three of the four instances, including both
sixth-order cases, and trails a Fourier network only on the smoothest
fourth-order amplitude.

\begin{table}[htbp]
\centering
\small
\caption{High-order real operators: relative \(L^2\) error \eqref{eq:relL2}.
``order'' is the operator order; the swept parameter is the order \(2m\)
(polyharmonic), the plate mode \(m\), or the KdV wavenumber \(k\). The plain
real \(\sinh\) baseline (omitted) exceeds \(0.85\) on every high-order row.}
\label{tab:res-ho}
\begin{tabular}{@{}llrrrr@{}}
\toprule
family & sweep & complex \(\sinh\) & Fourier & SIREN & MscaleDNN \\
\midrule
polyharmonic 1D & \(2m{=}2\)  & \textbf{4.15e-5} & 2.33e-4 & 2.13e-4 & 3.09e-4 \\
                & \(2m{=}4\)  & \textbf{2.03e-3} & 1.20e-2 & 3.15e-2 & 6.83e-2 \\
                & \(2m{=}6\)  & 1.42e-1 & 6.90e-1 & 9.66e-1 & \textbf{1.25e-1} \\
                & \(2m{=}8\)  & \textbf{2.14e-1} & 1.00e0  & 1.00e0  & 4.08e-1 \\
                & \(2m{=}10\) & \textbf{9.95e-1} & 1.00e0  & 1.01e0  & 9.98e-1 \\
\midrule
polyharmonic 2D & \(2m{=}2\)  & \textbf{4.23e-4} & 6.66e-4 & 7.83e-4 & 1.09e-2 \\
                & \(2m{=}4\)  & 3.96e-2 & \textbf{3.17e-2} & 9.83e-1 & 8.99e-1 \\
                & \(2m{=}6\)  & \textbf{7.67e-2} & 9.99e-1 & 9.94e-1 & 9.96e-1 \\
\midrule
plate (iso)     & \(m{=}1\)   & 4.62e-2 & \textbf{4.24e-2} & 9.88e-1 & 9.80e-1 \\
                & \(m{=}2\)   & \textbf{2.47e-1} & 2.76e-1 & 5.35e-1 & 1.00e0 \\
                & \(m{=}3\)   & \textbf{3.63e-1} & 6.72e-1 & 9.24e-1 & 1.00e0 \\
\midrule
plate (mixed)   & \(m{=}2\)   & \textbf{6.59e-1} & 9.78e-1 & 9.39e-1 & 1.00e0 \\
                & \(m{=}3\)   & \textbf{5.54e-1} & 9.93e-1 & 9.81e-1 & 1.00e0 \\
                & \(m{=}4\)   & \textbf{6.90e-1} & 9.97e-1 & 9.88e-1 & 1.00e0 \\
\midrule
beam (1D)       & \(m{=}1\)   & 4.36e-2 & \textbf{1.35e-2} & 1.01e0  & 9.78e-2 \\
                & \(m{=}2\)   & \textbf{1.29e-1} & 1.57e-1 & 6.11e-1 & 5.49e-1 \\
                & \(m{=}3\)   & 2.22e-1 & \textbf{1.16e-1} & 8.06e-1 & 1.68e0 \\
\midrule
linearized KdV  & \(k{=}2\)   & 3.65e-1 & \textbf{2.89e-1} & 4.54e-1 & 1.29e0 \\
                & \(k{=}3\)   & 5.68e-1 & 6.39e-1 & \textbf{6.05e-1} & 2.62e0 \\
                & \(k{=}4\)   & 8.70e-1 & 7.22e-1 & \textbf{6.55e-1} & 2.51e0 \\
                & \(k{=}5\)   & 9.48e-1 & 1.06e0  & \textbf{6.82e-1} & 4.03e0 \\
                & \(k{=}6\)   & 8.03e-1 & 1.10e0  & \textbf{7.37e-1} & 3.12e0 \\
\midrule
Cahn--Hilliard  & \(2m{=}4,a2\) & 4.45e-1 & \textbf{3.65e-1} & 5.55e-1 & 1.30e0 \\
                & \(2m{=}4,a3\) & \textbf{4.10e-1} & 5.90e-1 & 7.68e-1 & 1.51e0 \\
                & \(2m{=}6,a2\) & \textbf{5.60e-1} & 8.02e-1 & 9.74e-1 & 1.33e0 \\
                & \(2m{=}6,a3\) & \textbf{8.34e-1} & 9.11e-1 & 9.99e-1 & 1.57e0 \\
\bottomrule
\end{tabular}
\end{table}

% ---------------------------------------------------------------------------
\subsection{Genuinely complex-valued fields}\label{ssec:res-cplx}

For complex-valued targets the real baselines carry the field as a split
real--imaginary pair (an \(\mathbb{R}\)-valued PINN). Table~\ref{tab:res-cplx}
collects the cubic nonlinear Schr\"odinger soliton and the time-harmonic Maxwell
problem in a lossy medium. The complex \(\sinh\) network is the most accurate
method on every instance of both families: it leads the cubic NLS by a factor of
\(1.6\)--\(2.2\) over the best split-real network and reaches a relative error
below \(10^{-2}\) at \(k=1\), and on lossy Maxwell the split-real \(\tanh\) pair
fails outright (errors above \(1.5\)). This is the regime where a complex value
field is most natural, and the measured advantage is the most consistent in the
whole suite.

\begin{table}[htbp]
\centering
\small
\caption{Genuinely complex-valued problems: relative \(L^2\) error
\eqref{eq:relL2}. Real baselines carry a split real--imaginary pair.}
\label{tab:res-cplx}
\begin{tabular}{@{}llrrrr@{}}
\toprule
family & \(k\) / \(a\) & complex \(\sinh\) & SIREN & Fourier & split \(\tanh\) \\
\midrule
cubic NLS       & \(k{=}1\) & \textbf{7.56e-3} & 1.20e-2 & 1.82e-2 & 2.76e-1 \\
                & \(k{=}2\) & \textbf{3.26e-2} & 7.06e-2 & 7.90e-2 & 3.84e-1 \\
                & \(k{=}4\) & \textbf{4.13e-2} & 8.82e-2 & 8.95e-2 & 3.95e-1 \\
\midrule
Maxwell (lossy) & \(a{=}2\) & \textbf{2.31e-1} & 3.63e-1 & 3.94e-1 & 1.76e0 \\
                & \(a{=}4\) & \textbf{4.84e-1} & 7.62e-1 & 8.29e-1 & 1.82e0 \\
                & \(a{=}6\) & \textbf{6.22e-1} & 9.38e-1 & 1.02e0  & 1.55e0 \\
\bottomrule
\end{tabular}
\end{table}

% ---------------------------------------------------------------------------
\subsection{Summary and acceptance criterion}\label{ssec:res-summary}

Table~\ref{tab:res-summary} aggregates the contest by family. Counting an
instance as a win when the complex network attains the strictly smallest
relative \(L^2\) error, the complex \(\sinh\) ansatz is the best of the field on
\(32\) of the \(45\) instances of this section. The wins are not uniform: they
concentrate on the regimes the ansatz is designed for -- balanced high-order
operators (the polyharmonic, plate, and Cahn--Hilliard families) and genuinely
complex-valued fields (NLS and Maxwell, where the complex network wins every
instance) -- and on non-separable high-frequency targets (the chirp), while the
gains shrink or reverse on the first-order-in-time dispersive KdV, on the
one-dimensional beam and low-frequency plate modes, and on the low-frequency or
strongly directional Helmholtz problems.

We had set the stronger \emph{acceptance criterion} that the advantage factor
\eqref{eq:adv} grow monotonically with order or wavenumber. The data do not
support this in its strict form: although the per-family advantage is almost
always above one, it does not increase monotonically, and at the most extreme
order or wavenumber every method saturates near \(\mathcal{E}_{L^2}\approx 1\)
within the fixed budget, so the ratio is no longer informative. The honest
reading is therefore the weaker but robust one: \emph{at matched parameters and
a shared protocol, a complex value field is competitive-to-superior across the
oscillatory suite and clearly superior on balanced high-order and
complex-valued problems, but it is not a universal win and does not overcome the
budget-limited saturation of high-frequency PINNs.}

\begin{table}[htbp]
\centering
\small
\caption{Per-family summary. ``wins'' counts instances on which the complex
\(\sinh\) network attains the smallest relative \(L^2\) error; ``adv'' is the
range of the advantage factor \eqref{eq:adv} over the family's instances.}
\label{tab:res-summary}
\begin{tabular}{@{}llrl@{}}
\toprule
family & regime & wins & adv range \\
\midrule
Helmholtz (high \(k\))       & high-freq.\ 2nd order   & 5/5 & 1.07--1.55 \\
Helmholtz (anis.\ \(1,4\))   & directional 2nd order   & 0/1 & 0.57 \\
Helmholtz (var.\ \(\kappa\)) & heterogeneous 2nd order & 1/3 & 0.59--1.13 \\
radial chirp                 & non-separable 2nd order & 4/4 & 1.15--1.88 \\
\midrule
polyharmonic 1D              & order sweep \(2m\le10\)  & 4/5 & 0.88--5.94 \\
polyharmonic 2D              & order sweep \(2m\le6\)   & 2/3 & 0.80--12.96 \\
plate (isotropic)            & 4th order, modes        & 2/3 & 0.92--1.85 \\
plate (mixed modes)          & 4th order, anisotropic  & 3/3 & 1.42--1.77 \\
beam (1D)                    & 4th order, 1D           & 1/3 & 0.31--1.22 \\
linearized KdV               & 3rd order, dispersive   & 1/5 & 0.72--1.07 \\
Cahn--Hilliard               & 4th/6th, nonlinear      & 3/4 & 0.82--1.44 \\
\midrule
cubic NLS                    & complex-valued          & 3/3 & 1.59--2.16 \\
Maxwell (lossy)              & complex-valued          & 3/3 & 1.51--1.57 \\
\midrule
\textbf{total}               &                         & \textbf{32/45} & \\
\bottomrule
\end{tabular}
\end{table}
